\section{Information-theoretic proposals: the mass of information and ``infodynamics''}
\label{sec:information}

One reading of the simulation hypothesis holds that if the world is computed, information should
behave as a physical quantity with dynamics of its own. The most widely publicized proposals in this
vein are by M.~M.~Vopson: a \emph{mass--energy--information} (M/E/I) equivalence principle and a
\emph{second law of information dynamics} (``infodynamics''). The second law is presented as
evidence that ``appears to underpin'' the simulated-universe hypothesis~\cite{Vopson2023second}.
This section keeps three things apart:
(i)~the established thermodynamics of information,
(ii)~the conjectures built on it, and
(iii)~what would test those conjectures, and what the tests show so far.
Statements labelled ``our calculation'' are our own checks. The accompanying research notes give
the details needed to reproduce them.

\subsection{Established background: Landauer's principle}
\label{sec:information:landauer}

Landauer argued that logically irreversible operations are associated with physical irreversibility
and require a minimal heat generation ``typically of the order of $kT$ for each irreversible
function''~\cite{Landauer1961irreversibility}. His central calculation concerns resetting a thermal
ensemble of bits. This lowers their entropy by $k\ln 2$ per bit, so $0.6931\,kT$ per restored bit
must appear as heat ``to the surroundings''~\cite{Landauer1961irreversibility}.
Three further points in the same paper matter below.
\begin{itemize}
\item Devices such as a particle in a symmetric bistable well ``can hold information without
  dissipating energy''.
\item Resetting bits that are already known to be in the target state involves ``no entropy
  changes'' and ``no heat dissipation''.
\item When stored information thermalizes, the entropy of the information-bearing degrees of
  freedom ``can increase by $kN\log_e 2$''.
\end{itemize}
Bennett's formulation makes the bookkeeping explicit~\cite{Bennett2003notes}. An entropy decrease
of the information-bearing degrees of freedom must be compensated by an entropy increase elsewhere,
typically as ``energy imported into the computer, converted to heat, and dissipated into the
environment''. Erasing random data can be thermodynamically reversible, whereas erasing known data
cannot~\cite{Bennett2003notes}.

The bound has been tested in several systems:
\begin{itemize}
\item \emph{Colloidal bead, double-well optical trap.} The mean dissipated heat saturates at the
  Landauer bound for long erasure cycles~\cite{Berut2012experimental}.
\item \emph{Colloidal particle, feedback trap.} The extrapolated mean work for full erasure is
  $0.71\pm0.03\,kT$, compared with $\ln 2\approx0.693$, and $0.05\,kT$ in a control protocol that
  does not erase. The authors note that this work ``is dissipated into a surrounding heat
  bath''~\cite{Jun2014high}.
\item \emph{Nanomagnetic bits.} The dissipation is $(4.2\pm0.9)$\,zJ, i.e.\
  $(1.0\pm0.22)\,k_BT$ at 300\,K, within two standard deviations of
  $k_BT\ln2$~\cite{Hong2016experimental}.
\end{itemize}
Modern treatments place these results within stochastic thermodynamics~\cite{Parrondo2015thermodynamics}.
In that framework the nonequilibrium free energy is $F=E-TS$, with $S$ the Shannon (von Neumann)
entropy. It exceeds its equilibrium value by $T\,D[\rho\Vert\rho_{\rm eq}]$, the temperature times
the relative entropy of the actual state with respect to equilibrium~\cite{Esposito2011second}.
The principle's derivations have critics: Norton argues that they rest on incorrect assumptions and
neglect fluctuations~\cite{Norton2011waiting}. None of the proposals below turns on that dispute,
however. We take the Landauer bound, with the heat going to the environment, as the established
baseline.

\subsection{The mass--energy--information conjecture}
\label{sec:information:mei}

\paragraph{Claims.}
Vopson proposed that a stored bit has a rest mass
\begin{equation}
m_{\rm bit}=k_BT\ln2/c^2,
\end{equation}
which gives $3.19\times10^{-38}$\,kg at 300\,K. He proposed a test: a data storage device should be
heavier when full than when erased, by $2.5\times10^{-25}$\,kg for 1\,Tb~\cite{Vopson2019mass}.
The rationale is a syllogism: ``if $M=E$ and $E=I$, then $M=E=I$''. On this view the Landauer
energy of a bit ``condenses into its equivalent mass-energy when the information is stored at
equilibrium''~\cite{Vopson2022second}.
The same numerical value appeared earlier with a different meaning. Herrera attached
${\approx}3\times10^{-35}$\,g to one bit as the \emph{decrease} of a system's mass that accompanies
the energy dissipated when a bit changes~\cite{Herrera2014mass}.

Later papers extend the idea in two steps.
\begin{itemize}
\item \emph{Information per particle.} Vopson takes the Shannon entropy of the relative
  abundances of particle species in baryonic matter as the ``information content per particle''.
  This gives 1.288 bits for $\{p,e,n\}$ and 1.509 bits for $\{u,d,e\}$, and about $6\times10^{80}$
  bits in the observable universe~\cite{Vopson2021estimation}. He postulates that stable massive
  particles store this information ``about themselves''~\cite{Vopson2021estimation}.
\item \emph{Wider consequences.} Earlier estimates tied $\sim\!10^{94}$ bits to the missing dark
  matter; the 2021 paper itself calls this ``most likely overestimated''~\cite{Vopson2021estimation}.
  Growth of digital data was projected to exceed half of Earth's mass in $\sim500$
  years~\cite{Vopson2020information}.
\end{itemize}

\paragraph{The proposed test.}
The experimental protocol combines these ingredients~\cite{Vopson2022experimental}.
\begin{itemize}
\item \emph{Temperature-dependent mass.} A body at temperature $T$ should carry an information
  mass. For 1\,kg of copper changed by 100\,K the prediction is
  $\Delta m_{\rm inf}=3.33\times10^{-11}$\,kg. The paper states that the ``physical mass'' does
  not change with temperature.
\item \emph{Annihilation photons.} Because the electron's putative information mass is
  $\sim\!2\times10^{7}$ times smaller than its rest mass, the paper instead proposes to ``erase''
  it by $e^+e^-$ annihilation. In addition to the two 511\,keV photons, two photons of wavelength
  $\lambda=hc/(Ik_BT\ln2)$ should appear. This is $\sim$50\,$\mu$m at room temperature for
  $I=1.509$ bits and 25--75\,$\mu$m for 1--3 bits. We compute 45.9\,$\mu$m for $I=1.509$ at
  300\,K, i.e.\ 0.027\,eV per photon.
\item \emph{Apparatus.} A $^{22}$Na source, a W(100) moderator, a nm-thin Al target and
  temperature control. The predicted linear shift of $\lambda$ with $T$ is the main control.
\end{itemize}
The protocol states its assumptions openly. It assumes that the annihilating particles share the
target temperature, and it makes ``a strong assumption'' that the information energy leaves as IR
photons rather than, for example, with the gamma rays. Hence it is ``not totally guaranteed to
succeed'' even if the conjectures are correct~\cite{Vopson2022experimental}. A null result would
therefore not falsify M/E/I.

\paragraph{Status of tests.}
The proponent states that M/E/I ``still awaits an experimental confirmation''~\cite{Vopson2022second}
and ``has no empirical validation yet''~\cite{Vopson2025response}. We are not aware of any reported
execution of the annihilation protocol. The only direct weighings of storage media that we found
had an accuracy of $10^{-8}$\,kg. They recorded unexplained transient weight changes of order
1\,mg, which their authors did not attribute to a mass of information~\cite{Kish2013does}. That
accuracy is about 17 orders of magnitude short of the predicted $2.5\times10^{-25}$\,kg per
terabyte.

\paragraph{Critiques and assessment.}
Two published critiques address the conjecture directly. Burgin and Mikkilineni argue that only
physical \emph{representations} of information have mass, and that these can differ for the same
information~\cite{Burgin2022is}. Lairez argues on thermodynamic grounds~\cite{Lairez2024on}. At
constant temperature the internal energy of a system of independent entities is constant, so an
entropy change $T\Delta S$ is exchanged with the surroundings and cannot be localized as rest mass.
He adds that erasing the same bit twice, or erasing after making an independent copy, would expose
the inconsistency of attributing mass to logical information. A 2024 review of the Landauer bound
records these objections~\cite{Bormashenko2024landauer}. Lairez also disputes Landauer's principle
itself, but the mass argument does not depend on that dispute.

Our assessment agrees with Lairez's thermodynamic point and rests only on the standard results
above.
\begin{enumerate}
\item \emph{A stored bit in a symmetric memory carries no extra rest mass.} Landauer's symmetric
  memory stores a bit with no energy difference between its logical states. Kish and Granqvist
  make the same point for such memories: the system's ``energy and mass \dots\ remain the same''
  after writing~\cite{Kish2013does}. A bit known to be in one well therefore has the same mean
  energy $E$ as the thermalized bit. It differs only in nonequilibrium free energy, by
  $T D[\rho\Vert\rho_{\rm eq}]=k_BT\ln2$~\cite{Esposito2011second}. Since the mass of a body
  measures its energy content, $m=E_0/c^2$~\cite{Okun2008einstein}, this entropic free-energy
  excess implies no extra rest mass. The $k_BT\ln 2$ that appears on erasure is supplied as work
  by the eraser and leaves as heat~\cite{Jun2014high,Bennett2003notes}. It is not released from
  the bit's rest energy. For asymmetric memories a mass difference of either sign does exist, but
  it is set by the device's energy asymmetry~\cite{Kish2013does}, not by $k_BT\ln2$.
\item \emph{The copper prediction already conflicts with calorimetry (our calculation).} Take the
  conjecture literally, including $E=mc^2$ and energy conservation. Then the temperature-dependent
  $\Delta m_{\rm inf}$ of the protocol~\cite{Vopson2022experimental} must be supplied as energy
  when the body is heated. For copper this is an extra heat capacity of
  $1.509\times219.5\,k_B\ln2$ per atom, i.e.\ $1.9$\,kJ\,mol$^{-1}$\,K$^{-1}$. That is about 78
  times the measured $C_p\approx24.5$\,J\,mol$^{-1}$\,K$^{-1}$, which we obtain from the
  NIST-JANAF parameters~\cite{NIST2026copper}. Equivalently, $\Delta m_{\rm inf}c^2\approx3.0$\,MJ
  for the 1\,kg, 100\,K example, whereas the measured enthalpy increment between 300 and 400\,K is
  $3.9\times10^{4}$\,J. The ordinary relativistic mass increase of the heated block,
  $\Delta H/c^2\approx4\times10^{-13}$\,kg, is non-zero and about 80 times smaller than the
  claimed information mass.
\item \emph{The ``information per particle'' is a property of the mixture (our calculation).} The
  1.288 and 1.509 bits are entropies of the species mix in a population, not properties of a
  single particle. A pure-hydrogen population would give exactly 1\,bit. The predicted IR
  wavelength inherits this dependence.
\end{enumerate}

\subsection{The ``second law of infodynamics''}
\label{sec:information:infodynamics}

\paragraph{Claims.}
Vopson and Lepadatu define the entropy of ``information-bearing states'' as
\begin{equation}
S_{\rm inf}=N k_B\ln2\,H(X),
\end{equation}
where $N$ is the number of information states and $H(X)$ is the Shannon entropy of the symbol
frequencies in the stored string. For example, $H=0.991$ for the 88-bit ASCII encoding of
``INFORMATION''~\cite{Vopson2022second}. They assert $\partial S_{\rm inf}/\partial t\le0$ and give
two examples.
\begin{itemize}
\item \emph{Digital storage.} In a micromagnetic Monte Carlo simulation, written bits self-erase by
  thermal relaxation. $N$, and with it $S_{\rm inf}$, falls. The decrease is said to be
  compensated by a rise in physical entropy.
\item \emph{SARS-CoV-2 genomes.} Ten genomes of identical length (29{,}903 nt) were selected from
  NCBI because they ``displayed an incremental number of SNP mutations with time'' (0--49 over 22
  months). Their nucleotide-composition entropy falls from 1.9570243 to 1.9562614 bits, ``rather
  linearly'' in time.
\end{itemize}
The authors read the genomic trend as hinting at a deterministic, entropy-driven mutation process.

Later work extends the law to atomic physics (Hund's rules), to the prevalence of symmetry, and to
cosmology~\cite{Vopson2023second,Vopson2025is}. The cosmological derivation assumed $dQ=T\,dS=0$
for adiabatic expansion. After criticism that this holds only for reversible processes, Vopson
proposed a modified argument with added ``loss'' terms~\cite{Vopson2025on}. Finally, a toy model of
Planck-area cells storing 0 (empty) or 1 (occupied) is used to argue that gravitational clustering
lowers $H$. It derives $F=GMm/R^2$ by combining $\Delta S_{\rm inf}=k_B\ln2\,H$ per step
$\Delta r=\hbar/mc$ with $M=NHk_BT\ln2/c^2$ and $N\approx R^2/L_P^2$. On this basis gravity is
presented as ``data compression'' that ``supports the possibility of a simulated or computational
universe''~\cite{Vopson2025is}.

\paragraph{Independent analyses and critiques.}
\begin{itemize}
\item \emph{Replication with large datasets.} Ostrowski \textit{et al.}\ fitted composition-entropy
  trends to 8.8 million SARS-CoV-2, 1.1 million HIV and 1 million influenza records and found
  negative slopes throughout~\cite{Ostrowski2024entropy}. For SARS-CoV-2 the slope is
  $(-0.972\pm0.018)\times10^{-5}$ per mutation by least squares, with $\chi^2=7.6\times10^{7}$,
  and $(-1.146\pm0.421)\times10^{-5}$ by robust Bayesian regression. The robust slopes for HIV and
  influenza lie within $1.4\sigma$ of zero. The authors interpret the trend through nonequilibrium
  thermodynamics rather than a new law, and do not test a mutational-bias null model.
\item \emph{Reformulations.} Crecraft attributes the magnetic example to the ordinary second law.
  He argues that the statistical information entropy used is ``ill defined'' as the basis of a
  physical law, and notes that unbiased random mutations ``would be expected to increase'' it. He
  then proposes a substantially different ``thermocontextual'' law~\cite{Crecraft2024second}.
  Todd reinterprets the decreasing quantity as a relative entropy to equilibrium, whose decrease
  is compatible with the second law~\cite{Todd2026thermodynamic}.
\item \emph{Other applications.} Kubi\'nski \textit{et al.}\ report a U-shaped (eventually
  increasing) Shannon entropy in simulated nuclear fuel~\cite{Kubinski2024shannon}.
\item \emph{Public critique of the gravity paper.} Hossenfelder argued in a video that $H$ is
  unchanged by swapping 0s and 1s and is maximal at equal numbers, and that the paper ``shouldn't
  have been published''~\cite{Hossenfelder2025gravity}. Vopson replied that the video gives no
  explicit mathematical rebuttal and that fixing an encoding gives bit values
  meaning~\cite{Vopson2025response}. The reply and the cosmological note appeared as ``News and
  Views'' items in \emph{IPI Letters}, a journal edited by Vopson. The journal states that such
  items are screened at editorial level rather than going through the usual peer review.
\end{itemize}

\paragraph{Assessment (our analysis).}
\begin{enumerate}
\item \emph{Genome data.} We downloaded the ten accessions and reproduced every published entropy
  to seven decimals, using the plug-in estimator on A/C/G/U counts. The decrease is then a direct
  consequence of the known mutational spectrum.
  \begin{itemize}
  \item Moving one site from base $a$ to base $b$ changes $H$ by
    $\approx N^{-1}\log_2(p_a/p_b)$, which is $-2.7\times10^{-5}$\,bits for C$\to$U.
  \item SARS-CoV-2 is U-rich and C-poor (32.1\% vs 18.4\%), and 38--42\% of its early changes are
    C$\to$U, attributed to host APOBEC-like editing~\cite{Simmonds2020rampant}.
  \item In the most mutated genome of the table, 22 of 49 differences are C$\to$U, and 41 of 49
    move a site toward a more common base.
  \item A null model with unbiased substitutions predicts an \emph{increase} of
    $+3.9\times10^{-6}$ bits per substitution, compared with the observed $-1.6\times10^{-5}$.
  \end{itemize}
  The choice of estimator is not the issue here. At fixed $N$ the Miller--Madow correction is a
  constant offset~\cite{Paninski2003estimation}. What matters is the absence of a null model and
  the selection of sequences ``with an incremental number'' of SNPs. A monotone composition drift
  under a biased mutation process is expected on standard population-genetic grounds and does not
  indicate a new law.
\item \emph{Digital example.} For the magnetic medium, Landauer's accounting says that
  thermalization \emph{raises} the entropy of the information-bearing degrees of freedom by up to
  $Nk_B\ln2$~\cite{Landauer1961irreversibility}. What falls is the relative entropy the medium
  retains about the written record. For undriven relaxation, $\Delta I=-\Delta_iS\le0$ follows
  from Eq.~(12) of Ref.~\cite{Esposito2011second}. The ``second law of infodynamics'' in this
  example is therefore a restatement of the ordinary second law, as Crecraft and Todd also
  conclude.
\item \emph{Gravity model.} In the gravity model $H$ depends only on the counts $(N_0,N_1)$. It is
  unchanged when a particle moves through empty cells and changes only when particles merge. The
  $1/R^2$ law enters through the identifications $\Delta r=\hbar/mc$, $M=NHk_BT\ln2/c^2$ and
  $N\approx R^2/L_P^2$, not through a configuration-dependent entropy~\cite{Vopson2025is}.
  Moreover, it uses the M/E/I relation, whose difficulties are discussed above.
\end{enumerate}

\subsection{Summary}
\label{sec:information:summary}

Table~\ref{tab:information} summarizes the status of each claim. The empirically established core,
Landauer's bound with heat delivered to the environment, gives no support to a rest mass of stored
information. As formulated, the M/E/I conjecture conflicts with the standard energy bookkeeping of
symmetric memories. Its temperature-dependent version conflicts with calorimetry by nearly two
orders of magnitude. Its flagship test is untested, and it has a declared escape route if no
photons are seen. The infodynamics trends are reproducible arithmetic on well-defined data. Their
sign is explained by a known mutational bias in the genomic case and by the ordinary second law in
the digital case. None of these results, therefore, bears on whether the universe is
computed~\cite{Vopson2023second,Vopson2025is}. The infodynamics analyses are cheap to audit, and the
research notes outline a protocol for doing so.

\begin{table}[t]
\centering
\small
\caption{Status of the information-theoretic claims discussed in
Sec.~\ref{sec:information}.}
\label{tab:information}
\begin{tabular}{p{0.26\linewidth}p{0.2\linewidth}p{0.22\linewidth}p{0.22\linewidth}}
\hline
Claim & Status & Decisive test & Test status \\
\hline
Landauer bound, $\ge k_BT\ln2$ of heat per erased bit~\cite{Landauer1961irreversibility} &
Established; derivations debated~\cite{Norton2011waiting} &
Single-bit erasure calorimetry &
Confirmed within errors~\cite{Berut2012experimental,Jun2014high,Hong2016experimental} \\
Stored bit has rest mass $k_BT\ln2/c^2$~\cite{Vopson2019mass} &
Conjecture; contested~\cite{Lairez2024on,Burgin2022is} &
Weighing a full vs.\ erased memory ($\sim$$10^{-25}$\,kg/Tb) &
Not performed; best weighings are ${\sim}10^{-8}$\,kg~\cite{Kish2013does} \\
Temperature-dependent information mass of matter~\cite{Vopson2022experimental} &
Conjecture &
Mass or heat capacity vs.\ $T$ &
In conflict with calorimetry (our calculation) \\
$\sim$50\,$\mu$m photons from $e^+e^-$ annihilation~\cite{Vopson2022experimental} &
Prediction, with a stated escape clause &
Coincident far-IR detection with $T$-dependent $\lambda$ &
No experiment reported \\
Second law of infodynamics~\cite{Vopson2022second} &
Proposed law; reinterpreted~\cite{Crecraft2024second,Todd2026thermodynamic} &
Null models for mutation bias; entropy accounting &
Trends reproduced~\cite{Ostrowski2024entropy}; explained by mutational bias and the ordinary second law (our analysis) \\
Gravity as information compression~\cite{Vopson2025is} &
Conjecture; criticized informally~\cite{Hossenfelder2025gravity} &
None proposed &
--- \\
\hline
\end{tabular}
\end{table}
